% ============================================================================
\section{Research Gaps and Future Directions}
\label{sec:future_directions}
% ============================================================================

SPB operation, modulation, balancing, and modular control are
established, but the evidence needed for automotive architecture
selection remains incomplete. The key gaps are not additional
proof-of-concept control schemes but high-power validation and
system-level comparisons under common boundary conditions.

\subsection{High-Power Validation and Standardized Benchmarking}

Building on the laboratory-scale validation status noted in
Section~\ref{sec:evolution}, future
prototypes should target several-hundred-kilowatt scale
with realistic dc voltage, thermal, switching-overshoot, and EMI
constraints, among others,
benchmarked directly against a two- or three-level reference drive under
the same machine requirements, semiconductor generation, cooling
boundary, and mission. The scale of benefit that motivates this effort
is not merely hypothetical: a segmented, interleaved traction-drive
system operating at 55~kW has reported a 55--89\% reduction in dc-bus
capacitor ripple current relative to a conventional single-bridge
drive \cite{Su2012SegmentedTraction}, an order of magnitude above the
power level of most direct SPB studies and a useful automotive-scale
reference point for the interleaving benefit discussed in
Section~\ref{sec:modulation}.

Standardized reporting is also needed: published studies vary in power
level, cooling, filtering, and semiconductor generation, and
comparisons can shift materially with the chosen switching-energy
interpolation, temperature correction, and current-derating
assumptions. Reference operating points, efficiency maps, measured
parasitics, and open device/loss-model data would improve
reproducibility as semiconductor technology evolves.

A practical development path should also separate the questions that
can be answered at different hardware scales; Table~\ref{tab:validation_ladder}
proposes a validation ladder so that a low-power proof of one control
function is not read as evidence of automotive readiness.

\begin{table}[!t]
\centering
\caption{Suggested Validation Ladder for Automotive SPB Research}
\label{tab:validation_ladder}
\renewcommand{\arraystretch}{1.05}
\setlength{\tabcolsep}{3pt}
\begin{tabular}{p{1.55cm}p{2.55cm}p{3.35cm}}
\hline
\textbf{Stage} & \textbf{Primary purpose} & \textbf{Minimum evidence}\\
\hline
Cell / double-pulse & Device and commutation feasibility & Overshoot, switching energy, thermal impedance, protection\\
Multi-cell bench & Balancing and modulation & Cell-voltage spread, ripple, timing mismatch, transient stability\\
Dynamometer & Converter--machine interaction & Torque--speed map, efficiency, harmonics, CM/EMI, thermal limits\\
Mission / vehicle-relevant & Architecture value & Drive-cycle energy, cooling, fault derating, reliability and packaging\\
\hline
\end{tabular}
\end{table}

A community benchmark defining one passenger-car and one heavy-duty
traction case -- common dc-voltage range, power, torque--speed
envelope, coolant temperature, modulation limits, and EMC boundary,
with a VECTO-derived mission for heavy-duty work -- would let competing
architectures be compared under comparable electrical and thermal duty
without forcing a single reference machine. Reporting should extend the
same way: beyond peak efficiency, prototypes should state cell-voltage
spread, capacitor ripple current, common-mode current, and post-fault
capability, among other operating-point metrics, linking the
modulation, balancing, semiconductor, and reliability literatures into
transferable evidence.

\subsection{Cell Count, Reliability, and Cost Co-Design}

The optimum cell count $N$ remains the central unresolved design
question: as established in Section~\ref{sec:operating_principle},
the voltage-versus-component-count tradeoff shifts the optimum as
SiC, GaN, packaging, and cooling evolve
\cite{WBGTractionReview2025,Post800V2026}. A
common-platform comparison of a few-cell 650-V GaN SPB against a
many-cell low-voltage GaN SPB, using real commercial devices, would be
valuable; the GaN comparison in \cite{Rohner2024GaN} is a useful
starting point but should be extended to automotive power with current
parallelization, machine voltage/speed constraints, dc-link
capacitance, cooling, and mission energy.

Converter, machine, and thermal models, among others, should
be co-optimized rather than designed sequentially. Existing phase-count
studies show machine design changes with SPB cell number
\cite{Bringezu2020SPB}; the missing step is a complete automotive
co-design using vehicle-level energy, volume, mass, cost, and
reliability as a joint objective.

Reliability is part of that objective: future studies should combine
semiconductor, capacitor, sensor, gate-driver, communication, and
winding faults (Section~\ref{sec:fault_reconfiguration};
\cite{Nikouie2016Fault,Verkroost2024Reconfiguration}) in one
state-based, Markov-style framework as used for related fault-tolerant
IMMDs \cite{Swanke2022MarkovReliability}, keeping \emph{component
failure}, \emph{module isolation}, and \emph{loss of vehicle function}
distinct.
An $N$-cell drive may survive with $N-1$ active modules, but its
torque--speed envelope, module thermal loading, and dc-voltage margin
all change, so reliability data should pair with a post-fault
capability model -- not a series reliability block diagram -- to
compare availability, limp-home power, and mission completion
probability against efficiency.

Cost and sustainability assessment should use the same $N$-dependent
framing: increasing $N$ can reduce per-switch die cost but tends to
increase isolated gate supplies, sensing channels, and assembly
operations, among others, almost linearly, while a
repeated converter--winding module can gain from common tooling and
automated assembly -- a balance likely to differ between high-volume
passenger vehicles and lower-volume heavy-duty systems. The same
accounting should separate production, operation, maintenance, and
end-of-life impact over one vehicle mission and service life, since
modular repair and longer service life can offset the added production
impact of extra components.

Overall, SPB research must move from demonstrating that the
architecture operates toward demonstrating where it is better: whether
it gives a measurable vehicle-level advantage over mature two- and
three-level architectures once machine design, thermal management,
EMC, reliability, manufacturing, and cost are included.
